% !TEX program = xelatex
% Статья о конструктивной форме леммы Гензеля; сборка выполняется XeLaTeX.

\documentclass{misisarticle}

% Коды намеренно не указаны: их выбирают по правилам конкретного сборника.
\grnti{}
\asjc{}
\oecd{}

\title{Лемма Гензеля на практике: от решения по модулю $p$ к $p$-адическому корню}

% Замените поля на сведения автора перед подачей статьи.
\misisauthor{В.~С.~Немаев}{кафедра инженерной кибернетики НИТУ «МИСиС»}{Москва, Россия}{m2300915@edu.misis.ru}
\misisauthor{В.~В.~Шихеева}{кафедра инженерной кибернетики НИТУ «МИСиС»}{Москва, Россия}{valyashik@gmail.com}


\abstract{Рассматривается конструктивная форма леммы Гензеля для многочленов с целыми коэффициентами. Цель работы состоит в том, чтобы связать поиск корня сравнения по модулю простого числа с вычислением приближений к корню в кольце $p$-адических целых чисел. Сформулировано достаточное условие единственного подъёма простого корня. Выведена рекуррентная формула, позволяющая на каждом шаге находить одну новую $p$-адическую цифру корня решением линейного сравнения. Алгоритм применён к многочлену $x^2-2$ при $p=7$. Пример показывает, что арифметика по модулю последовательных степеней $p$ даёт точные приближения без операций с вещественными числами.}

\keywords{лемма Гензеля, $p$-адические числа, сравнения, подъём корня, теория чисел}

\begin{document}

\maketitle

\section{Введение}

Сравнения по модулю простого числа часто служат первым и наиболее доступным этапом исследования диофантовых уравнений. Однако решение сравнения $f(x)\equiv0\pmod p$ само по себе не отвечает на вопрос о разрешимости по более высоким степеням $p$ и о существовании корня в $p$-адических числах. Лемма Гензеля даёт точный переход от конечной арифметики вычетов к $p$-адическому пределу~\cite{hensel1908,serre1973}.

Практическая ценность результата состоит в его конструктивности. Если корень по модулю $p$ простой, то его можно последовательно уточнять по модулям $p^2,p^3,\ldots$. На каждом шаге требуется вычислить значение многочлена и решить одно линейное сравнение по модулю $p$. Такой подход используется при работе с локальными полями, в алгоритмах компьютерной алгебры и при локальном анализе диофантовых задач~\cite{koblitz1984,gouvea1997}.

Цель работы --- описать этот переход в форме алгоритма, доказать корректность одного шага уточнения и разобрать полный вычислительный пример.

\section{Анализ работ}

Идея подъёма корня принадлежит К.~Гензелю, который ввёл $p$-адические числа как средство переноса разложений и решений сравнений на бесконечную глубину~\cite{hensel1908}. В классическом курсе Ж.-П.~Серра лемма формулируется как критерий существования и единственности корня в $\mathbb Z_p$ при ненулевой производной по модулю $p$~\cite{serre1973}.

В учебных изложениях обычно подчёркивается аналогия с методом Ньютона. В вещественном анализе метод Ньютона является приближённым и требует контроля сходимости. В $p$-адической ситуации при простом корне каждый шаг леммы Гензеля точно фиксирует очередной разряд числа по основанию $p$~\cite{gouvea1997}. Это делает метод особенно удобным для ручных вычислений и программной реализации.

Обобщённые варианты леммы допускают кратные корни, но их условия зависят от $p$-адических порядков значений $f(a)$ и $f'(a)$~\cite{koblitz1984}. В данной работе рассматривается простая форма: она содержит основной механизм подъёма и не требует дополнительных оценок.

\section{Постановка задачи}

Пусть $p$ --- простое число, $f(x)\in\mathbb Z[x]$, а $a_1\in\mathbb Z$ удовлетворяет условиям
\begin{equation}
  f(a_1)\equiv0\pmod p,
  \qquad
  f'(a_1)\not\equiv0\pmod p.
  \label{eq:simple-root}
\end{equation}
Необходимо построить последовательность целых чисел $a_n$, для которой
\begin{equation}
  f(a_n)\equiv0\pmod {p^n},
  \qquad
  a_{n+1}\equiv a_n\pmod {p^n}.
  \label{eq:lift-goal}
\end{equation}

Последовательность сравнений во второй части~\eqref{eq:lift-goal} означает, что классы $a_n\pmod {p^n}$ согласованы. Поэтому они задают единственное число $\alpha\in\mathbb Z_p$. Требуется показать, что $f(\alpha)=0$, и получить правило вычисления $a_{n+1}$ по известному $a_n$.

\section{Метод подъёма простого корня}

\subsection{Один шаг уточнения}

Пусть уже найдено $a_n$ такое, что $p^n\mid f(a_n)$. Следующее приближение имеет вид
\begin{equation}
  a_{n+1}=a_n+t p^n,
  \qquad t\in\{0,1,\ldots,p-1\}.
  \label{eq:next-approximation}
\end{equation}
Для многочлена с целыми коэффициентами разложение по степеням приращения даёт
\begin{equation}
  f(a_n+t p^n)\equiv f(a_n)+t p^n f'(a_n)
  \pmod {p^{n+1}}.
  \label{eq:taylor}
\end{equation}
Так как $f(a_n)=p^n q_n$ для некоторого целого $q_n$, из~\eqref{eq:taylor} следует условие на новую цифру:
\begin{equation}
  q_n+t f'(a_n)\equiv0\pmod p.
  \label{eq:digit-equation}
\end{equation}

Производная $f'(a_n)$ сравнима с $f'(a_1)$ по модулю $p$, поскольку $a_n\equiv a_1\pmod p$. По условию~\eqref{eq:simple-root} она обратима в $\mathbb F_p$. Поэтому у сравнения~\eqref{eq:digit-equation} есть единственное решение
\begin{equation}
  t\equiv-q_n\bigl(f'(a_n)\bigr)^{-1}\pmod p.
  \label{eq:digit-formula}
\end{equation}
Именно это значение подставляется в~\eqref{eq:next-approximation}.

\subsection{Корректность и единственность}

Формула~\eqref{eq:digit-formula} доказывает существование единственного $a_{n+1}$ по модулю $p^{n+1}$, продолжающего заданный класс $a_n\pmod {p^n}$. Индукция по $n$ строит всю последовательность. Для любого фиксированного $N$ все члены, начиная с $a_N$, совпадают по модулю $p^N$; следовательно, последовательность фундаментальна в $\mathbb Z_p$ и имеет предел $\alpha$.

Многочлен непрерывен в $p$-адической метрике. Переходя к пределу в делимости $p^n\mid f(a_n)$ при всех $n$, получаем $f(\alpha)=0$. Если другой корень $\beta$ удовлетворяет $\beta\equiv a_1\pmod p$, то на каждом шаге он обязан выбирать ту же единственную цифру $t$. Значит, $\beta=\alpha$. Таким образом, простой корень по модулю $p$ поднимается единственным образом.

\subsection{Алгоритм}

Для вычисления приближения по модулю $p^N$ достаточно повторить следующий шаг $N-1$ раз:
\begin{enumerate}
  \item вычислить $q_n=f(a_n)/p^n$;
  \item найти обратный к $f'(a_n)$ элемент по модулю $p$;
  \item вычислить $t$ по формуле~\eqref{eq:digit-formula};
  \item положить $a_{n+1}=a_n+t p^n$.
\end{enumerate}

Все вычисления обратного элемента происходят в поле из $p$ элементов. Поэтому при фиксированном $p$ один шаг использует только целочисленные операции и не требует приближённой арифметики.

\section{Вычислительный пример}

Рассмотрим $f(x)=x^2-2$ и $p=7$. Число $a_1=3$ является корнем по модулю $7$, поскольку $3^2-2=7$. Кроме того, $f'(x)=2x$ и $f'(3)=6\not\equiv0\pmod7$, так что применима простая форма леммы.

На первом шаге $q_1=f(3)/7=1$. Обратный к $f'(3)=6$ по модулю $7$ равен $6$, поэтому
\[
  t\equiv-1\cdot6\equiv1\pmod7,
  \qquad a_2=3+1\cdot7=10.
\]
Действительно, $10^2-2=98=2\cdot7^2$. На втором шаге $q_2=98/49=2$ и $f'(10)=20\equiv6\pmod7$. Получаем $t\equiv-2\cdot6\equiv2\pmod7$, откуда $a_3=10+2\cdot49=108$. Продолжая те же вычисления, получаем шесть последовательных уточнений и приближение $a_7$ по модулю $7^7$.

\begin{table}[htbp]
  \caption{Последовательные приближения к $\sqrt{2}$ в $\mathbb Z_7$}
  \label{tab:hensel-example}
  \centering
  \setlength{\tabcolsep}{1.5pt}
  \begin{tabular}{c c c c >{\raggedright\arraybackslash}p{4.2cm}}
    \toprule
    $n$ & $a_n$ & $q_n=f(a_n)/7^n$ & $t_n$ & Разложение $a_n$ \\
    \toprule
    1 & 3 & 1 & 1 & $3$ \\
    2 & 10 & 2 & 2 & $3+1\cdot7^1$ \\
    3 & 108 & 34 & 6 & $3+1\cdot7^1+2\cdot7^2$ \\
    4 & 2166 & 1954 & 1 & $\begin{gathered}3+1\cdot7^1+2\cdot7^2\\{}+6\cdot7^3\end{gathered}$ \\
    5 & 4567 & 1241 & 2 & $\begin{gathered}3+1\cdot7^1+2\cdot7^2+6\cdot7^3\\{}+1\cdot7^4\end{gathered}$ \\
    6 & 38181 & 12391 & 1 & $\begin{gathered}3+1\cdot7^1+2\cdot7^2+6\cdot7^3\\{}+1\cdot7^4+2\cdot7^5\end{gathered}$ \\
    7 & 155830 & 29486 & --- & $\begin{gathered}3+1\cdot7^1+2\cdot7^2+6\cdot7^3\\{}+1\cdot7^4+2\cdot7^5+1\cdot7^6\end{gathered}$ \\
    \bottomrule
  \end{tabular}
\end{table}

Здесь $a_{n+1}=a_n+t_n7^n$. В частности, $155830^2-2=29486\cdot7^7$, поэтому $a_7=155830$ является решением сравнения $x^2\equiv2\pmod {7^7}$. Таблица~\ref{tab:hensel-example} показывает, что каждое новое значение сохраняет все предыдущие $7$-адические цифры и добавляет одну новую. Предел последовательности является единственным $7$-адическим корнем $x^2-2$, сравнимым с $3$ по модулю $7$. Второй корень, сравнимый с $4\equiv-3\pmod7$, поднимается независимо и равен $-\alpha$.

\begin{thebibliography}{9}

\bibitem{hensel1908} Hensel K. Theorie der algebraischen Zahlen. Leipzig: B. G. Teubner, 1908.

\bibitem{serre1973} Serre J.-P. A Course in Arithmetic. New York: Springer-Verlag, 1973. 115 p.

\bibitem{koblitz1984} Koblitz N. $p$-adic Numbers, $p$-adic Analysis, and Zeta-Functions. 2nd ed. New York: Springer-Verlag, 1984. 151 p.

\bibitem{gouvea1997} Gouv\^ea F. Q. $p$-adic Numbers: An Introduction. 2nd ed. Berlin: Springer-Verlag, 1997. 298 p.

\end{thebibliography}

\end{document}
